\documentclass[12pt]{article}
\usepackage[T1]{fontenc}
\usepackage[french]{babel}
\usepackage{amsmath}
\usepackage{amssymb}
\usepackage{graphicx}
\usepackage{xcolor}
\usepackage{lmodern}

\title{\textbf{Bayesian Knowledge Tracing (BKT) - Version Complète}}
\author{D'après Corbett \& Anderson (1994, 2005)}
\date{\today}

\begin{document}

\maketitle

\tableofcontents

\newpage

\section{Introduction}

\subsection{Fondements théoriques}

Ce document présente l'algorithme BKT complet tel que décrit dans les papiers fondateurs :
\begin{itemize}
    \item \textbf{Corbett, K. R., \& Anderson, J. R. (1994).} ``Knowledge Tracing: Modeling the Acquisition of Procedural Knowledge''
    \item \textbf{Corbett, K. R., \& Anderson, J. R. (2005).} ``Knowledge Tracing: A Generalized Low-Order Markov Model for Representing Learning''
\end{itemize}

La principale différence avec les implémentations simplifiées :
\begin{itemize}
    \item Tous les 4 paramètres ($p_0, p_T, p_G, p_S$) sont estimés
    \item Utilisation d'un algorithme EM complet avec forward-backward pass
    \item Inférence probabiliste stricte sur les états cachés
\end{itemize}

\newpage

\section{Modèle probabiliste}

\subsection{Définition formelle}

BKT est un \textbf{chaîne de Markov cachée} (Hidden Markov Model - HMM) à temps discret.

\subsubsection{État caché}

L'état de connaissance de l'étudiant à l'étape $t$ est une variable binaire :

$$L_t \in \{0, 1\}$$

où :
\begin{itemize}
    \item $L_t = 0$ : L'étudiant ne connaît pas la compétence
    \item $L_t = 1$ : L'étudiant connaît la compétence
\end{itemize}

\subsubsection{Observation}

La réponse observée à l'étape $t$ est :

$$C_t \in \{0, 1\}$$

où :
\begin{itemize}
    \item $C_t = 1$ : Réponse correcte
    \item $C_t = 0$ : Réponse incorrecte
\end{itemize}

\subsection{Hypothèse d'indépendance conditionnelle}

L'état futur dépend seulement de l'état présent (propriété de Markov) :

$$P(L_{t+1} \mid L_1, L_2, \ldots, L_t) = P(L_{t+1} \mid L_t)$$

\section{Les quatre paramètres}

\subsection{1. Prior : $p_0$}

$$p_0 = P(L_0 = 1)$$

\textbf{Signification :} Probabilité a priori que l'étudiant connaisse la compétence avant toute tentative.

\textbf{Plage typique :} $[0.1, 0.5]$

\subsection{2. Learn (Transition) : $p_T$}

$$p_T = P(L_{t+1} = 1 \mid L_t = 0)$$

\textbf{Signification :} Probabilité que l'étudiant transite de l'état ne connaît pas vers connaît après une tentative.

\textbf{Plage typique :} $[0.1, 0.3]$

\subsection{3. Guess (Devinette) : $p_G$}

$$p_G = P(C_t = 1 \mid L_t = 0)$$

\textbf{Signification :} Probabilité d'une réponse correcte malgré l'absence de connaissance.

\textbf{Plage typique :} $[0.05, 0.2]$

\subsection{4. Slip (Glissement) : $p_S$}

$$p_S = P(C_t = 0 \mid L_t = 1)$$

\textbf{Signification :} Probabilité d'une réponse incorrecte malgré la connaissance.

\textbf{Plage typique :} $[0.05, 0.2]$

\section{Dynamique du modèle}

\subsection{Transition d'état}

La probabilité de transition entre états :

$$P(L_{t+1} = 1 \mid L_t) = \begin{cases}
        p_T & \text{si } L_t = 0 \\
        1   & \text{si } L_t = 1
    \end{cases}$$

Ou de manière compacte :

$$P(L_{t+1} = 1 \mid L_t) = L_t + (1-L_t)p_T$$

\subsection{Observation générée}

La probabilité d'observer une réponse donnée l'état :

$$P(C_t = 1 \mid L_t) = \begin{cases}
        p_G     & \text{si } L_t = 0 \\
        1 - p_S & \text{si } L_t = 1
    \end{cases}$$

Ou de manière compacte :

$$P(C_t = 1 \mid L_t) = L_t(1-p_S) + (1-L_t)p_G$$

\newpage

\section{Algorithme EM complet}

\subsection{Vue d'ensemble}

L'algorithme EM alterne entre :

\begin{enumerate}
    \item \textbf{E-step} : Calculer distribution a posteriori des états cachés
    \item \textbf{M-step} : Réestimer les paramètres
\end{enumerate}

\subsection{E-STEP : Forward-Backward Pass}

\subsubsection{Forward Pass}

Pour chaque étudiant, calculer :

$$\ell_t^{\text{fwd}} = P(L_t = 1 \mid C_{1:t})$$

\textbf{Initialisation :}
$$\ell_0^{\text{fwd}} = p_0$$

\textbf{Prédiction avant observation :}

$$\ell_t^{\text{pred}} = \ell_{t-1}^{\text{fwd}} + (1-\ell_{t-1}^{\text{fwd}})p_T$$

\textbf{Prédiction de l'observation :}

$$P(C_t = 1 \mid C_{1:t-1}) = \ell_t^{\text{pred}}(1-p_S) + (1-\ell_t^{\text{pred}})p_G$$

\textbf{Mise à jour si } $C_t = 1$ :

$$\ell_t^{\text{fwd}} = \frac{(1-p_S) \ell_t^{\text{pred}}}{\ell_t^{\text{pred}}(1-p_S) + (1-\ell_t^{\text{pred}})p_G}$$

\textbf{Mise à jour si } $C_t = 0$ :

$$\ell_t^{\text{fwd}} = \frac{p_S \ell_t^{\text{pred}}}{\ell_t^{\text{pred}} p_S + (1-\ell_t^{\text{pred}})(1-p_G)}$$

\subsubsection{Backward Pass}

Calculer la probabilité lissée :

$$\ell_t^{\text{bwd}} = P(L_t = 1 \mid C_{1:T})$$

\textbf{Initialisation :}

$$\ell_T^{\text{bwd}} = \ell_T^{\text{fwd}}$$

\textbf{Récurrence (pour } $t = T-1, \ldots, 0$) :

$$\ell_t^{\text{bwd}} = \ell_t^{\text{fwd}} \cdot \frac{\ell_{t+1}^{\text{bwd}}}{\ell_t^{\text{pred}}} + (1-\ell_t^{\text{fwd}}) \cdot \frac{1-\ell_{t+1}^{\text{bwd}}}{1-\ell_t^{\text{pred}}}$$

\section{M-STEP : Réestimation des paramètres}

\subsection{Expected counts}

Accumuler pendant E-step :

$$\mathbb{E}[n_{L_t=0}] = \sum_{i,t} (1 - \ell_t^{\text{bwd}})$$

$$\mathbb{E}[n_{L_t=1}] = \sum_{i,t} \ell_t^{\text{bwd}}$$

Observations correctes sans connaissance :

$$\mathbb{E}[n_{C_t=1, L_t=0}] = \sum_{i,t : C_t=1} (1-\ell_t^{\text{bwd}}) \cdot p_G$$

Observations incorrectes avec connaissance :

$$\mathbb{E}[n_{C_t=0, L_t=1}] = \sum_{i,t : C_t=0} \ell_t^{\text{bwd}} \cdot p_S$$

\subsection{Formules de réestimation}

\textbf{Prior :}

$$p_0^{\text{new}} = \frac{1}{N} \sum_{i=1}^{N} \ell_0^{\text{bwd}, i}$$

\textbf{Learn :}

$$p_T^{\text{new}} = \frac{\text{Transitions apprises}}{\mathbb{E}[n_{L_t=0}]}$$

\textbf{Guess :}

$$p_G^{\text{new}} = \frac{\mathbb{E}[n_{C_t=1, L_t=0}]}{\mathbb{E}[n_{C_t=1, L_t=0}] + \mathbb{E}[n_{C_t=0, L_t=0}]}$$

\textbf{Slip :}

$$p_S^{\text{new}} = \frac{\mathbb{E}[n_{C_t=0, L_t=1}]}{\mathbb{E}[n_{C_t=1, L_t=1}] + \mathbb{E}[n_{C_t=0, L_t=1}]}$$

\subsection{Clamping}

$$p_i^{\text{new}} = \max(\min(p_i^{\text{new}}, 0.99), 0.01)$$

\section{Critère de convergence}

L'algorithme converge quand :

$$\sum_{i \in \{p_0, p_T, p_G, p_S\}} |p_i^{\text{new}} - p_i^{\text{old}}| < 10^{-4}$$

\newpage

\section{Pseudocode de l'algorithme}

\subsection{Entraînement EM}

Pour chaque compétence $s$ :

\begin{enumerate}
    \item Initialiser : $p_0 = 0.3, p_T = 0.15, p_G = 0.1, p_S = 0.1$

    \item Pour chaque itération $k$ :
          \begin{enumerate}
              \item \textbf{E-STEP}
                    \begin{itemize}
                        \item Pour chaque étudiant $i$ :
                              \begin{itemize}
                                  \item Forward pass : calculer $\ell_t^{\text{fwd}}$
                                  \item Backward pass : calculer $\ell_t^{\text{bwd}}$
                                  \item Accumuler expected counts
                              \end{itemize}
                    \end{itemize}

              \item \textbf{M-STEP}
                    \begin{itemize}
                        \item Réestimer $p_0, p_T, p_G, p_S$
                        \item Clamp à $[0.01, 0.99]$
                    \end{itemize}

              \item Vérifier convergence
          \end{enumerate}

    \item Stocker paramètres : $\text{params}[s] = [p_0, p_T, p_G, p_S]$
\end{enumerate}

\subsection{Prédiction}

Pour chaque (étudiant, compétence) :

\begin{enumerate}
    \item $\ell_0 \leftarrow p_0$
    \item Pour $t = 1 \ldots T$ :
          \begin{enumerate}
              \item $\ell_t^{\text{pred}} \leftarrow \ell_{t-1} + (1-\ell_{t-1})p_T$
              \item $\text{prob} \leftarrow \ell_t^{\text{pred}}(1-p_S) + (1-\ell_t^{\text{pred}})p_G$
              \item Stocker prédiction
              \item Mettre à jour selon $C_t$
          \end{enumerate}
\end{enumerate}

\section{Différences clés}

\subsection{Forward-Backward vs Forward only}

\begin{table}[ht]
    \centering
    \begin{tabular}{|l|c|c|}
        \hline
        \textbf{Aspect} & \textbf{Forward only} & \textbf{Forward-Backward} \\
        \hline
        Obs futures     & Non                   & Oui                       \\
        Probabilités    & Filtrées              & Lissées                   \\
        Précision EM    & Approx                & Exacte                    \\
        Complexité      & $O(NT)$               & $O(NT)$                   \\
        \hline
    \end{tabular}
\end{table}

\section{Considérations pratiques}

\subsection{Données}

\begin{itemize}
    \item Triées chronologiquement par étudiant
    \item Groupées par compétence
    \item Pas de valeurs manquantes
\end{itemize}

\subsection{Initialisation}

\begin{itemize}
    \item $p_0 = 0.3$
    \item $p_T = 0.15$
    \item $p_G = 0.1$
    \item $p_S = 0.1$
\end{itemize}

\subsection{Itérations}

\begin{itemize}
    \item Petit dataset : 3-5 itérations
    \item Dataset moyen : 5-10 itérations
    \item Grand dataset : 10-20 itérations
\end{itemize}

\subsection{Stabilité}

\begin{itemize}
    \item Clamp : $p_i \in [0.01, 0.99]$
    \item Epsilon : $\hat{C}_t = \max(\hat{C}_t, 10^{-10})$
\end{itemize}

\newpage

\section{Exemple numérique}

Étudiant avec 5 observations.

Paramètres : $p_0=0.3, p_T=0.15, p_G=0.1, p_S=0.1$

Observations : $C = [1, 1, 0, 1, 0]$

\subsection{Forward Pass}

\begin{align}
    t=0: \quad \ell_0^{\text{fwd}}  & = 0.3                                       \\
    t=1: \quad \ell_1^{\text{pred}} & = 0.405, \quad \ell_1^{\text{fwd}} = 0.870  \\
    t=2: \quad \ell_2^{\text{pred}} & = 0.8895, \quad \ell_2^{\text{fwd}} = 0.985 \\
    t=3: \quad \ell_3^{\text{pred}} & = 0.9873, \quad \ell_3^{\text{fwd}} = 0.878 \\
    t=4: \quad \ell_4^{\text{pred}} & = 0.878, \quad \ell_4^{\text{fwd}} = 0.51
\end{align}

L'étudiant apprend rapidement (2 bonnes réponses), mais une mauvaise réponse crée de l'incertitude.

\section{Conclusion}

L'algorithme EM complet pour BKT :

\begin{enumerate}
    \item Estime les 4 paramètres
    \item Forward-backward pour raffiner
    \item Convergence vers maximum vraisemblance
    \item Stable numériquement
\end{enumerate}

Approche rigoureuse des papiers fondateurs.

\end{document}